\section{Conclusion}
In this paper we have introduced a general category theoretic foundation for epistemic syntax, which cements uncertainty calculus into a common framework on a higher level than what is usually done. 

Our approach allowed us to naturally construct mathematical incantations of epistemic philosophical positions and concepts, and to have a natural way to compare them. 

We have also introduced categorical versions of known examples of epistemic syntaxes. This shows that it is versatile, and allows for a new perspective on existing examples. We also described which philosophical positions and concepts hold for each of these examples. 

We then introduced a structured approach for comparing epistemic syntaxes, together with an understanding of how such changes affects the logic and epistemic quality. 

Finally, we introduced a formal method for assigning epistemic uncertainty to a given structure -- in the form of a general category. As category theory is very general, and allows for the modeling of a lot of structures, this gives a robust and mathematically sound framework for discussing the uncertainties assigned to it, and how it combines along features of the structure. This method also allowed us to construct a mathematically rigorous process for changing the epistemic syntax assigned to a system. 

There are bound to be other positions and concepts that can be modeled in this framework, and we hope that this paper can be a starting point to a broader discussion and exploration on categorical modeling of epistemic ideas. 

